%% toc: P3 31 31.9
\chapter{Self-Extending Models and Indiscernibles of Arbitrary Length}
\label{sec:p03-w03-p03-030-self-extending}

Some theories come with a way of building, inside any of their models, a
canonical proper elementary extension of that model.  Such a theory can be
iterated: extend, extend again, take unions at limit stages, and continue
for as many steps as wanted.  A long iteration produces models whose order
type is controlled by the iteration, for instance models in which every
proper initial segment is small while the model itself is large.  It also
produces long well-ordered sequences of indiscernible elements, one new
element for each step.  This chapter presents that method, called the method
of \emph{self-extending models}, and its best-known application, the
MacDowell--Specker--Shelah theorem on proper elementary end-extensions,
proved here for models of the induction scheme and the least-number
principle in a countable vocabulary.

The presentation follows Section~6 of Schmerl's chapter on transfer theorems
in the 1985 Springer edition of \emph{Model-Theoretic Logics}, pages
202--208, read in the open-access copy of that edition
\cite[pp.~202--208]{Schmerl1985TransferTheorems}.  Locators below refer to
that edition.  Where the source states a result without proof, or sketches a
proof, the gap is stated at the place it occurs, and only arguments that are
written out in full in the annex are counted as proved here.

\emph{Inputs.}  A first-order vocabulary \(L\) (usually containing a binary
symbol \(<\)), a theory in the extended language \(\mathscr L(Q)\) obtained
by adding one unary quantifier \(Q\), and a model of it.  \emph{Outputs.}
Canonical elementary extensions; chains of them of any ordinal length;
order-indiscernible sequences of that length; strongly \(\lambda\)-like
models for regular \(\lambda\) exceeding the size of the language; and, for
theories with the induction scheme and the least-number principle in a
countable vocabulary, proper elementary end-extensions, in particular for
models of Peano arithmetic; and \(\lambda\)-like models for every set of
fewer than \(\lambda\) first-order sentences whose finite subsets have models
of order type \(\omega\).

The finitary first-order indiscernibles of
\cref{def:p03-order-indiscernibles} and the Ehrenfeucht--Mostowski
construction of \cref{thm:p03-ehrenfeucht-mostowski} obtain indiscernibles
from Ramsey's theorem and compactness.  The method here obtains them instead
from an ultrafilter that the theory itself describes, and it gives
additional control over the order in which the new elements sit.  The
remaining transfer and two-cardinal material of Schmerl's chapter, which uses
the iteration below as one of its proof methods, is treated with the other
sections of this chapter.

\section{The logic \texorpdfstring{\(\mathscr L(Q)\)}{L(Q)} with an ultrafilter reading}
\label{sec:p03-w03-p03-030-logic}

\begin{hbdefinition}{def:p03-w03-p03-030-weak-structures}{Structures for \(\mathscr L(Q)\) and reduced structures}
Let \(L\) be a first-order vocabulary.  The formulas of \(\mathscr L(Q)\) are
built like first-order formulas over \(L\) with one more formation rule: if
\(\phi\) is a formula and \(x\) a variable, then \(Qx\,\phi\) is a formula in
which \(x\) is bound.  A structure for \(\mathscr L(Q)\) is a pair
\((\mathfrak A,q)\), where \(\mathfrak A\) is an \(L\)-structure with domain
\(A\) and \(q\) is a collection of subsets of \(A\).  Satisfaction is defined
as in first-order logic, with the extra clause
\[
 (\mathfrak A,q)\entails Qx\,\phi(x)
 \quad\text{iff}\quad
 \{a\in A:(\mathfrak A,q)\entails\phi(a)\}\in q ,
\]
for formulas with parameters from \(A\).  A subset of \(A\) is
\emph{definable} when it has the form
\(\{a\in A:(\mathfrak A,q)\entails\phi(a,\bar b)\}\) for an
\(\mathscr L(Q)\)-formula \(\phi\) and parameters \(\bar b\) from \(A\).
The structure is \emph{reduced} if every member of \(q\) is definable
\cite[p.~203]{Schmerl1985TransferTheorems}.

Replacing \(q\) by its definable members changes no satisfaction fact,
because the clause for \(Q\) only asks about definable sets; so every
structure can be replaced by a reduced one with the same theory.  A theory in
\(\mathscr L(Q)\) is \emph{consistent} here when it has a model of this kind.
\index{self-extending model!structure for L(Q)}
\end{hbdefinition}

The quantifier \(Q\) is not interpreted by a fixed cardinality condition.
Each structure carries its own family \(q\), and the theories below force
\(q\) to behave like a nonprincipal ultrafilter on the definable sets.

\begin{hbdefinition}{def:p03-w03-p03-030-self-extending-theory}{Self-extending theory}
A consistent theory \(T\) of \(\mathscr L(Q)\) is \emph{self-extending}
\cite[Definition~6.1.1, pp.~202--203]{Schmerl1985TransferTheorems} when it
has the following two properties.
\begin{enumerate}
\item \emph{\(T\) is a Skolem theory.}  For every formula
  \(\phi(x_0,\ldots,x_{n-1},y)\) of \(\mathscr L(Q)\) there is an
  \(L\)-term \(t(x_0,\ldots,x_{n-1})\) such that
  \[
   \forall\bar x\,\bigl(\exists y\,\phi(\bar x,y)\rightarrow
                         \phi(\bar x,t(\bar x))\bigr)
  \]
  is in \(T\).
\item \emph{\(Q\) behaves as a nonprincipal ultrafilter.}  The universal
  closures of all formulas of the following four forms are in \(T\):
  \begin{gather*}
   Qx\,\phi(x)\leftrightarrow\neg Qx\,\neg\phi(x),\qquad
   \forall y\,\neg Qx\,(x=y),\\
   \forall x\,(\phi(x)\rightarrow\psi(x))\rightarrow
     (Qx\,\phi(x)\rightarrow Qx\,\psi(x)),\\
   Qx\,\phi(x)\wedge Qx\,\psi(x)\rightarrow Qx\,(\phi(x)\wedge\psi(x)).
  \end{gather*}
\end{enumerate}
The source writes the Skolem term with a script letter; it is written
\(t\) here.
\index{self-extending theory}
\end{hbdefinition}

\begin{hbconvention}{def:p03-w03-p03-030-scheme-reading}{How the displayed schemes are read}
In the third ultrafilter scheme the source displays the antecedent as
\(\phi(x)\rightarrow\psi(x)\) inside a universal closure
\cite[p.~203]{Schmerl1985TransferTheorems}.  Taken literally, the closure
would quantify \(x\) over the whole implication, which is not the intended
monotonicity law; the quantifier \(\forall x\) is therefore written on the
antecedent above.

Two further schemes are used below.  The \emph{regressive scheme} of
\cite[Theorem~6.1.6, p.~204]{Schmerl1985TransferTheorems} consists of the
universal closures of
\[
 Qx\,\exists y\,(\phi(x,y)\wedge y<x)\rightarrow\exists y\,Qx\,\phi(x,y).
 \tag{R}
\]
The \emph{bounded scheme} of
\cite[Theorems~6.1.3 and 6.2.2, pp.~204 and 206]{Schmerl1985TransferTheorems}
is displayed there as
\(\forall z\,Qx\,\exists y(\phi(x,y)\wedge y<z)\rightarrow\exists y\,Qx\,\phi(x,y)\).
It is read here with \(\forall z\) governing the whole implication:
\[
 \forall z\,\bigl(Qx\,\exists y\,(\phi(x,y)\wedge y<z)\rightarrow
                  \exists y\,Qx\,\phi(x,y)\bigr).
 \tag{B}
\]
If \(\forall z\) governed only the antecedent, then in any linear order with a
least element the antecedent would fail at that element and the sentence
would say nothing.  Reading (B) is the one under which the scheme yields
end-extensions (\cref{thm:p03-w03-p03-030-end-extension}).
\end{hbconvention}

\begin{hblemma}{thm:p03-w03-p03-030-ultrafilter-facts}{proved=annex:p03-w03-p03-030-ultrafilter-facts}{The definable sets in \(q\) form a nonprincipal ultrafilter}
Let \((\mathfrak A,q)\) satisfy the four ultrafilter schemes of
\cref{def:p03-w03-p03-030-self-extending-theory}, with \(A\ne\varnothing\).
Let \(\mathcal D\) be the Boolean algebra of definable subsets of \(A\).
\begin{enumerate}
\item \(q\cap\mathcal D\) is an ultrafilter of \(\mathcal D\): it contains
  \(A\) and not \(\varnothing\), it is closed under finite intersections and
  under definable supersets, and for each \(X\in\mathcal D\) exactly one of
  \(X\) and \(A\setminus X\) is in \(q\).
\item No singleton is in \(q\); in particular \(A\) is infinite.
\item \((\mathfrak A,q)\entails Qx\,\phi\rightarrow\exists x\,\phi\) for every
  formula with parameters, and \((\mathfrak A,q)\entails
  Qx\,\chi\leftrightarrow\chi\) when \(x\) is not free in \(\chi\).
\end{enumerate}
\end{hblemma}

Because the family \(q\) is arbitrary, \(\mathscr L(Q)\) with this reading
is first-order logic in disguise: a quantifier \(Qx\,\phi\) can be replaced
by a new relation symbol, provided the new symbols are made to respect equality
of the sets they describe.  This gives compactness without any condition on
the size of the vocabulary.

\begin{hblemma}{thm:p03-w03-p03-030-compactness}{proved=annex:p03-w03-p03-030-compactness}{Compactness for structures for \(\mathscr L(Q)\)}
Let \(L\) be a first-order vocabulary of any cardinality and \(\Gamma\) a set
of \(\mathscr L(Q)\)-sentences over \(L\).  If every finite subset of
\(\Gamma\) has a model \((\mathfrak A,q)\) in the sense of
\cref{def:p03-w03-p03-030-weak-structures}, then \(\Gamma\) has such a model.
When \(L\) is uncountable the proof uses the axiom of choice, through the
first-order compactness theorem.
\end{hblemma}

\section{The canonical extension}
\label{sec:p03-w03-p03-030-canonical}

\begin{hbdefinition}{def:p03-w03-p03-030-canonical-extension}{Definable functions modulo \(q\) and the canonical extension}
Let \((\mathfrak A,q)\) be a reduced model of a self-extending theory \(T\).
A function \(f:A\to A\) is \emph{definable} when its graph is a definable
subset of \(A^2\).  For definable \(f_0,\ldots,f_{n-1}\) and an
\(\mathscr L(Q)\)-formula \(\phi(y_0,\ldots,y_{n-1})\), the expression
\(Qx\,\phi(f_0(x),\ldots,f_{n-1}(x))\) abbreviates the formula
\[
 Qx\,\exists y_0\cdots\exists y_{n-1}\Bigl(\bigwedge_{i<n}\gamma_i(x,y_i)\wedge
 \phi(\bar y)\Bigr),
\]
where \(\gamma_i\) defines the graph of \(f_i\).  It says that the set of
\(a\in A\) with \((\mathfrak A,q)\entails\phi(f_0(a),\ldots,f_{n-1}(a))\) is
in \(q\).

Two definable functions \(f,g\) are \emph{equal modulo \(q\)} when
\((\mathfrak A,q)\entails Qx\,(f(x)=g(x))\).  Let \(B\) be the set of classes
\([f]\) of definable functions under this relation
\cite[p.~203]{Schmerl1985TransferTheorems}.  The \emph{canonical extension}
\((\mathfrak B,r)\) of \((\mathfrak A,q)\) has domain \(B\); for each
relation symbol \(R\), function symbol \(F\) and constant symbol \(c\) of
\(L\),
\[
 R^{\mathfrak B}([f_0],\ldots,[f_{k-1}])
   \iff (\mathfrak A,q)\entails Qx\,R(f_0(x),\ldots,f_{k-1}(x)),
\]
\(F^{\mathfrak B}([f_0],\ldots,[f_{k-1}])=[x\mapsto
F^{\mathfrak A}(f_0(x),\ldots,f_{k-1}(x))]\) and \(c^{\mathfrak B}=[x\mapsto
c^{\mathfrak A}]\).  The family \(r\) consists of all sets
\[
 \{[g]\in B:(\mathfrak A,q)\entails Qx\,\phi(f_0(x),\ldots,f_{n-1}(x),g(x))\}
\]
for which \((\mathfrak A,q)\entails Qx\,Qy\,\phi(f_0(x),\ldots,f_{n-1}(x),y)\)
\cite[p.~203]{Schmerl1985TransferTheorems}.  Each \(a\in A\) is identified
with the class of the constant function \(c_a\) with value \(a\), and the
class of the identity function is written \(\iota\).
\index{self-extending model!canonical extension}
\end{hbdefinition}

The construction is an ultrapower of \(\mathfrak A\) by \(q\) restricted to
definable functions.  Restricting to definable functions is what keeps the
extension inside the reach of the theory: the Skolem terms supply the
witnesses that a full ultrapower would take from arbitrary functions.

\begin{hbtheorem}{thm:p03-w03-p03-030-canonical-extension}{proved=annex:p03-w03-p03-030-canonical-extension}{The canonical extension and its \L o\'s equivalence}
Let \((\mathfrak A,q)\) be a reduced model of a self-extending theory \(T\)
in \(\mathscr L(Q)\), and let \((\mathfrak B,r)\) be its canonical extension.
\begin{enumerate}
\item The interpretations and the family \(r\) are well defined: they do not
  depend on the chosen representatives.
\item For every \(\mathscr L(Q)\)-formula \(\phi(y_0,\ldots,y_{n-1})\) without
  parameters and all definable \(f_0,\ldots,f_{n-1}\),
  \[
   (\mathfrak B,r)\entails\phi([f_0],\ldots,[f_{n-1}])
   \quad\text{iff}\quad
   (\mathfrak A,q)\entails Qx\,\phi(f_0(x),\ldots,f_{n-1}(x)).
  \]
\item \((\mathfrak B,r)\) is reduced, and it is the only reduced structure
  with domain \(B\) that satisfies the equivalence in (2).
\item The map \(a\mapsto[c_a]\) is an elementary embedding of
  \((\mathfrak A,q)\) into \((\mathfrak B,r)\) for all
  \(\mathscr L(Q)\)-formulas; in particular \((\mathfrak B,r)\entails T\).
\item \(\abs A\le\abs B\le\abs A+\abs L+\aleph_0\).
\end{enumerate}
\end{hbtheorem}

Parts (1)--(4) are the facts asserted in
\cite[p.~203]{Schmerl1985TransferTheorems}; part (5) is the cardinality
bound used in the iteration on
\cite[p.~205]{Schmerl1985TransferTheorems}.  The proof uses the Skolem
property at two places: for the existential quantifier and for the
quantifier \(Q\).  Without it, the extension could fail to be elementary.

\begin{hbcorollary}{thm:p03-w03-p03-030-large-sets}{proved=annex:p03-w03-p03-030-large-sets}{Large sets become larger}
In the situation of \cref{thm:p03-w03-p03-030-canonical-extension}, for every
\(\mathscr L(Q)\)-formula \(\phi(x,\bar y)\) and every tuple \(\bar a\) from
\(A\),
\[
 (\mathfrak B,r)\entails Qx\,\phi(x,\bar a)\rightarrow\phi(\iota,\bar a),
\]
and \(\iota\notin A\).  So every canonical extension is proper.
\end{hbcorollary}

This is the source's statement that ``large sets become larger''
\cite[p.~203]{Schmerl1985TransferTheorems}: a set that is large in the old
model acquires the new element \(\iota\).

\begin{hblemma}{thm:p03-w03-p03-030-hull}{proved=annex:p03-w03-p03-030-hull}{Small models of a self-extending theory}
Every self-extending theory \(T\) in \(\mathscr L(Q)\) over the vocabulary
\(L\) has a reduced model of cardinality at most \(\abs L+\aleph_0\).
\end{hblemma}

\begin{hblemma}{thm:p03-w03-p03-030-chains}{proved=annex:p03-w03-p03-030-chains}{Unions of elementary chains of reduced structures}
Let \(\delta\) be a limit ordinal and \(((\mathfrak A_\nu,q_\nu))_{\nu<\delta}\)
a sequence of reduced \(\mathscr L(Q)\)-structures that satisfy the
ultrafilter schemes, with \((\mathfrak A_\mu,q_\mu)\preccurlyeq
(\mathfrak A_\nu,q_\nu)\) for \(\mu<\nu<\delta\), elementarity being for all
\(\mathscr L(Q)\)-formulas.  Let \(\mathfrak A\) be the union of the
\(L\)-structures \(\mathfrak A_\nu\).  Put \(X\in q\) iff there are
\(\nu<\delta\), a formula \(\phi(x,\bar y)\) and \(\bar b\) from \(A_\nu\)
such that \((\mathfrak A_\nu,q_\nu)\entails Qx\,\phi(x,\bar b)\) and
\(X=\{a\in A:(\mathfrak A_\mu,q_\mu)\entails\phi(a,\bar b)\text{ for some
}\mu\ge\nu\text{ with }a\in A_\mu\}\).
Then \((\mathfrak A,q)\) is reduced and \((\mathfrak A_\nu,q_\nu)\preccurlyeq
(\mathfrak A,q)\) for every \(\nu<\delta\).

Consequently, if \((\mathfrak A_0,q_0)\) is a reduced model of a
self-extending theory and \((\mathfrak A_\nu,q_\nu)_{\nu\le\delta}\) is built
by taking canonical extensions at successor stages and such unions at limit
stages, then \((\mathfrak A_\mu,q_\mu)\preccurlyeq(\mathfrak A_\nu,q_\nu)\)
whenever \(\mu\le\nu\le\delta\), and every \((\mathfrak A_\nu,q_\nu)\) is a
reduced model of the same theory.
\end{hblemma}

The source takes ``the union of the previously constructed structures'' at
limit stages \cite[p.~205]{Schmerl1985TransferTheorems}; the family \(q\) in
the lemma is the one that makes this union an elementary extension of each
earlier stage.  The first-order part of the lemma is the elementary chain
theorem, \cref{thm:elementary-chain}.

\section{Iterating the extension: cardinal-like models and indiscernibles}
\label{sec:p03-w03-p03-030-iteration}

\begin{hbdefinition}{def:p03-w03-p03-030-cardinal-like}{\(\kappa\)-like, \((\kappa,X)\)-like and strongly \(\kappa\)-like orders}
Let \(\kappa\) be an infinite cardinal.  A linear order \((A,<)\) is
\emph{\(\kappa\)-like} when \(\abs A=\kappa\) and
\(\abs{\{b\in A:b<a\}}<\kappa\) for every \(a\in A\).  The source states the
initial-segment condition as the sentence \(\forall x\,\neg Qy\,(y<x)\) in
the \(\kappa\)-interpretation of \(Q\)
\cite[Definition~1.2.1, p.~179]{Schmerl1985TransferTheorems}.

Let \(\kappa\) be regular and uncountable and \(X\subseteq\kappa\).  A
\(\kappa\)-like \((A,<)\) is \emph{\((\kappa,X)\)-like} if there is an
increasing function \(e:X\to A\) such that whenever \(\alpha\in X\) and
\(\alpha=\sup\{\nu\in X:\nu<\alpha\}\), then
\(e(\alpha)=\sup\{e(\nu):\nu\in X\cap\alpha\}\), the supremum being the least
upper bound in \((A,<)\)
\cite[Definition~6.1.4, p.~204]{Schmerl1985TransferTheorems}.  The condition
is applied to nonzero \(\alpha\) only.  At \(\alpha=0\) the supremum of the
empty family would be the least element of \(A\), and \(\kappa\)-like orders
need not have one.  The order is
\emph{strongly \(\kappa\)-like} when it is \((\kappa,\kappa)\)-like
\cite[p.~204]{Schmerl1985TransferTheorems}; in the source's first
formulation, it is \(\kappa\)-like and contains a closed unbounded subset
\cite[Definition~5.2.1, p.~200]{Schmerl1985TransferTheorems}.  A structure
\((A,<,\ldots)\) has one of these properties when its order does.
\index{cardinal-like order!strongly}
\end{hbdefinition}

For example, the ordinal \(\kappa\) itself is strongly \(\kappa\)-like, with
\(e\) the identity.  The order \(\kappa\times\mathbb Q\), ordered
lexicographically, is \(\kappa\)-like when \(\kappa\) is uncountable: each
initial segment lies inside \((\alpha+1)\times\mathbb Q\) for some
\(\alpha<\kappa\), a set of size \(\abs\alpha+\aleph_0<\kappa\).  It is not
strongly \(\kappa\)-like.  An increasing \(e:\kappa\to\kappa\times\mathbb Q\)
puts at most countably many values in each copy \(\{\gamma\}\times\mathbb Q\),
so its first coordinates are unbounded in \(\kappa\).  Choose
\(\nu_0<\nu_1<\cdots\) whose first coordinates strictly increase, with
supremum \(\gamma\), and let \(\beta=\sup_k\nu_k\).  The values \(e(\nu)\)
for \(\nu<\beta\) then have first coordinates cofinal in the limit ordinal
\(\gamma\), and they have no least upper bound, because
\(\{\gamma\}\times\mathbb Q\) has no least element.  So continuity fails at
\(\beta\).

\begin{hblemma}{thm:p03-w03-p03-030-end-extension}{proved=annex:p03-w03-p03-030-end-extension}{End-extensions and least new elements}
Let \(T\) be a self-extending theory in which \(<\) is a linear order, let
\((\mathfrak A,q)\) be a reduced model of \(T\), and let \((\mathfrak B,r)\)
be its canonical extension.
\begin{enumerate}
\item If \(T\) contains the bounded scheme (B), then \(\mathfrak B\) is an
  end-extension of \(\mathfrak A\): every \(b\in B\) with \(b<a\) for some
  \(a\in A\) belongs to \(A\).
\item If \(T\) contains the regressive scheme (R) and the sentence
  \(\forall y\,Qx\,(y<x)\), then (B) holds in every model of \(T\), and
  \(\iota\) is the least element of \(B\setminus A\).
\end{enumerate}
\end{hblemma}

The source asserts that the sentences supplied by its Theorem~6.1.6, that is,
the regressive scheme, make every canonical extension an end-extension with a
least new element \cite[pp.~204--205]{Schmerl1985TransferTheorems}.  The
proof of part (2) uses, in addition, the sentence
\(\forall y\,Qx\,(y<x)\), which says that \(q\) contains every final segment.
That sentence is not among the displayed sentences of the source's theorem,
and it is listed as a hypothesis here.  The example in
\cref{def:p03-w03-p03-030-example-n} shows that part (1) can hold while a
least new element fails to exist.

\begin{hbtheorem}{thm:p03-w03-p03-030-strongly-like}{proved=annex:p03-w03-p03-030-strongly-like}{Strongly \(\lambda\)-like models from a self-extending theory}
Let \(T'\) be a self-extending theory in \(\mathscr L(Q)\) over \(L\), in
which \(<\) is a linear order, containing the regressive scheme (R) and the
sentence \(\forall y\,Qx\,(y<x)\).  Let \(\lambda\) be a regular cardinal with
\(\lambda>\abs L+\aleph_0\).  Start with a reduced model
\((\mathfrak A_0,q_0)\) of \(T'\) with \(\abs{A_0}<\lambda\).  Let
\((\mathfrak A_{\nu+1},q_{\nu+1})\) be the canonical extension of
\((\mathfrak A_\nu,q_\nu)\), take unions as in
\cref{thm:p03-w03-p03-030-chains} at limit stages, and let \(a_\nu\) be the
least element of \(A_{\nu+1}\setminus A_\nu\).  Then
\begin{enumerate}
\item \(A_\nu=\{x\in A_\lambda:x<a_\nu\}\) for every \(\nu<\lambda\);
\item \(\nu\mapsto a_\nu\) is increasing and continuous: for limit
  \(\alpha<\lambda\), \(a_\alpha\) is the least upper bound of
  \(\{a_\nu:\nu<\alpha\}\);
\item the \(L\)-reduct of \(\mathfrak A_\lambda\) is a strongly \(\lambda\)-like
  model of every first-order sentence of \(T'\).
\end{enumerate}
\end{hbtheorem}

This is the iteration described in the source
\cite[p.~205]{Schmerl1985TransferTheorems}, with its cardinal hypotheses: the
length of the chain is the regular cardinal \(\lambda\), and \(\lambda\) must
exceed the size of the language so that each stage stays below \(\lambda\).
The model required at stage \(0\) exists by \cref{thm:p03-w03-p03-030-hull}.

\begin{hbtheorem}{thm:p03-w03-p03-030-indiscernibles}{proved=annex:p03-w03-p03-030-indiscernibles}{The new elements form an indiscernible sequence of any prescribed length}
Let \((\mathfrak A_0,q_0)\) be a reduced model of a self-extending theory, let
\(\delta\ge1\) be an ordinal, and let \((\mathfrak A_\nu,q_\nu)_{\nu\le\delta}\)
be built by canonical extensions at successor stages and unions at limit
stages.  Let \(a_\nu\in A_{\nu+1}\) be the element \(\iota\) of the canonical
extension of \((\mathfrak A_\nu,q_\nu)\).  Then for every
\(\mathscr L(Q)\)-formula \(\phi(x_1,\ldots,x_n,\bar y)\), every \(\bar b\)
from \(A_0\) and all \(\nu_1<\cdots<\nu_n<\delta\),
\begin{multline*}
 (\mathfrak A_\delta,q_\delta)\entails
     \phi(a_{\nu_1},\ldots,a_{\nu_n},\bar b)\\
 \text{iff}\quad
 (\mathfrak A_0,q_0)\entails Qx_1\,Qx_2\cdots Qx_n\,\phi(x_1,\ldots,x_n,\bar b).
\end{multline*}
In particular the elements \(a_\nu\) \((\nu<\delta)\) are distinct, and the
sequence \((a_\nu)_{\nu<\delta}\) is order indiscernible over \(A_0\) in the
sense of \cref{def:p03-order-indiscernibles}, for all \(\mathscr L(Q)\)-formulas
and hence for all first-order formulas.
\end{hbtheorem}

The source does not state this indiscernibility; it is a consequence of the
\L o\'s equivalence and is proved in the annex.  Two features distinguish
these indiscernibles from those of \cref{thm:p03-ehrenfeucht-mostowski}.
Their index set is an ordinal of any length, chosen in advance, and the type
of every increasing tuple is computed from the single model
\((\mathfrak A_0,q_0)\) by prefixing quantifiers \(Q\).  Under the
hypotheses of \cref{thm:p03-w03-p03-030-strongly-like} the indiscernibles are
exactly the closed unbounded sequence \((a_\nu)_{\nu<\lambda}\) of that
theorem.

\begin{hblemma}{thm:p03-w03-p03-030-definable-ultrapower-n}{proved=annex:p03-w03-p03-030-definable-ultrapower-n}{A self-extending expansion of the natural numbers}
Let \(\mathcal U\) be a nonprincipal ultrafilter on \(\mathbb N\).  Let
\(L_0=\{<,+\}\), and let \(L_{k+1}\) add to \(L_k\) a function symbol
\(f_\chi\) for each \(L_k(Q)\)-formula \(\chi(\bar z,y)\), interpreted in
\(\mathbb N\) by the least \(y\) with \(\chi(\bar z,y)\), or by \(0\) when
there is none, the formula being evaluated in the structure with family
\(\mathcal U\).  Let \(L_\omega=\bigcup_kL_k\), let \(\mathfrak N\) be the
resulting \(L_\omega\)-structure on \(\mathbb N\), and let \(q\) be the set of
definable members of \(\mathcal U\).
\begin{enumerate}
\item The \(\mathscr L(Q)\)-theory \(T\) of \((\mathfrak N,q)\) is
  self-extending, and \(\forall y\,Qx\,(y<x)\in T\).
\item The canonical extension of \((\mathfrak N,q)\) is an end-extension.
\item The regressive scheme (R) fails in \((\mathfrak N,q)\), and the
  canonical extension has no least new element.
\end{enumerate}
\end{hblemma}

\begin{hbexample}{def:p03-w03-p03-030-example-n}{The natural numbers: an end-extension without a least new element}
In \cref{thm:p03-w03-p03-030-definable-ultrapower-n} a nonprincipal
ultrafilter on \(\mathbb N\) exists by the ultrafilter lemma
(\cref{thm:p01-ultrafilter-lemma}) applied to the filter of cofinite sets.
The canonical extension \((\mathfrak B,r)\) of \((\mathfrak N,q)\) is the
ultrapower of \(\mathfrak N\) by \(\mathcal U\) restricted to the definable
functions.  The new element \(\iota\) is the class of the identity, and
\(\iota>n\) for every \(n\in\mathbb N\) because cofinite sets are large.

Halving shows what the regressive scheme controls.  The function
\(h(x)=\lfloor x/2\rfloor\) is defined by \(y+y=x\vee y+y+1=x\).  We have
\(h(x)<x\) for \(x\ge1\), so \([h]<\iota\), and \([h]\) is new because
\(h\) takes each value at most twice.  Repeating with \(h\circ h\),
\(h\circ h\circ h\), and so on gives an infinite descending sequence of new
elements below \(\iota\).  So the iteration of
\cref{thm:p03-w03-p03-030-strongly-like} cannot be carried out with this
theory: there is no least new element to serve as \(a_\nu\), and the closed
unbounded set of that theorem is not available.  The extension is still an
end-extension, because a definable function bounded on a large set takes
only finitely many values there and is constant on a large subset.
\end{hbexample}

\section{The MacDowell--Specker--Shelah theorem}
\label{sec:p03-w03-p03-030-mss}

MacDowell and Specker showed that every model of Peano arithmetic has a
proper elementary end-extension.  Shelah showed that the argument needs only
an induction scheme for the order
\cite[p.~206]{Schmerl1985TransferTheorems}.  The source proves the result by
building a self-extending theory.  The construction searches a tree of
definable equivalence classes for a canonical branch of large classes, and
reads off the ultrafilter from that branch.

\begin{hbdefinition}{def:p03-w03-p03-030-induction-scheme}{Induction scheme, least-number principle and tuple orders}
Let \(L\) contain the binary symbol \(<\).  The \emph{induction scheme}
\cite[p.~206]{Schmerl1985TransferTheorems} is the sentence ``\(<\) is a
linear order with a first but no last element'' together with the universal
closures of all formulas
\[
 [\exists x\,\phi(x)\wedge\forall x\,\exists y\,(\phi(x)\rightarrow
   \phi(y)\wedge x<y)]\rightarrow\forall x\,\exists y\,(\phi(y)\wedge x<y).
\]
It says that a nonempty definable set with no largest element is cofinal.
The \emph{least-number principle} is the set of universal closures of
\(\exists x\,\phi(x)\rightarrow\exists x\,(\phi(x)\wedge\forall
y\,(y<x\rightarrow\neg\phi(y)))\).  The first element is written \(0\).

The tuple orders are defined by recursion
\cite[p.~207]{Schmerl1985TransferTheorems}: \(<_1\) is \(<\), and for
\((n+1)\)-tuples, \(\bar x<_{n+1}\bar y\) is the formula
\begin{multline*}
 \max(x_0,\ldots,x_n)<\max(y_0,\ldots,y_n)\vee
 \bigl[\max(x_0,\ldots,x_n)=\max(y_0,\ldots,y_n)\\
 \wedge\bigl(x_n<y_n\vee(x_n=y_n\wedge(x_0,\ldots,x_{n-1})<_n
 (y_0,\ldots,y_{n-1}))\bigr)\bigr].
\end{multline*}
For an element \(z\), \(z^{n+1}\) is the \((n+1)\)-tuple \((z,\ldots,z)\).
\index{induction scheme}\index{least-number principle}
\end{hbdefinition}

\begin{hblemma}{thm:p03-w03-p03-030-tuple-orders}{proved=annex:p03-w03-p03-030-tuple-orders}{Arithmetic of the tuple orders}
Let \(M\) satisfy the induction scheme and the least-number principle for
all formulas with parameters, and let \(n\ge1\).
\begin{enumerate}
\item \(<_n\) is a linear order on \(M^n\) with least element \(0^n\).
\item Every nonempty definable set of \(n\)-tuples has a \(<_n\)-least
  element.
\item Every nonempty definable set of \(n\)-tuples with a \(<_n\)-upper bound
  has a \(<_n\)-greatest element.
\item Every \(n\)-tuple has an immediate \(<_n\)-successor, and every
  \(n\)-tuple other than \(0^n\) has an immediate \(<_n\)-predecessor.
\item If \(\max(\bar x)<z\), then \(\bar x<_nz^n\); and \(z<z'\) implies
  \(z^n<_nz'^n\).
\item For every definable function \(F:M\to M\) and every \(m\in M\), the set
  \(\{F(z):z\le m\}\) has an upper bound.
\end{enumerate}
\end{hblemma}

Part (6) is a collection principle for definable functions; it is what the
proof below uses where the source speaks of ``boundedly'' many classes
\cite[p.~207]{Schmerl1985TransferTheorems}.

\begin{hbmethod}{def:p03-w03-p03-030-tree-search}{The tree of equivalence classes and its leftmost large branch}
Fix a formula \(\phi(\bar x,y)\), with \(\bar x\) an \((n+1)\)-tuple of
variables, and a formula \(\theta(y)\) defining an unbounded set.  The source
uses the following formulas \cite[pp.~207--208]{Schmerl1985TransferTheorems}:
\begin{align*}
 E(\bar x,y,z)&\defeq\forall\bar w<_{n+1}\bar x\,
               (\phi(\bar w,y)\leftrightarrow\phi(\bar w,z)),\\
 L(\bar x,y,z)&\defeq\exists\bar w<_{n+1}\bar x\,
     (\neg\phi(\bar w,y)\wedge\phi(\bar w,z)\wedge E(\bar w,y,z)),\\
 S(\bar x,y)&\defeq\forall w\,\exists z>w\,(\theta(z)\wedge E(\bar x,y,z))\\
   &\qquad\wedge\forall v\,\bigl(L(\bar x,v,y)\rightarrow
   \exists w\,\forall z>w\,(\theta(z)\rightarrow\neg E(\bar x,v,z))\bigr).
\end{align*}
For a fixed \emph{rank} \(\bar x\), \(E\) is an equivalence relation, and
it becomes finer as the rank increases; so the classes form a tree.  A class
is \emph{large} when it contains unboundedly many elements satisfying
\(\theta\).  \(L\) orders the classes of one rank lexicographically, by the
first parameter tuple at which they differ, and \(S(\bar x,y)\) says that
\(y\) lies in the first large class of rank \(\bar x\).  The selected classes
form a branch, and \(\sigma(\bar x)\defeq\forall\bar w\,\exists
y\,(S(\bar w,y)\wedge\phi(\bar x,y))\) reads off the value of \(\phi(\bar x,\cdot)\)
along that branch.

The source writes \(\theta_n\) where the construction requires
\(\theta_{n-1}\), and writes \(y'\) for the variable bound as \(v\); both are
read as in the display above.  Every step is justified in
\cref{thm:p03-w03-p03-030-tree}.
\end{hbmethod}

\begin{hblemma}{thm:p03-w03-p03-030-tree}{proved=annex:p03-w03-p03-030-tree}{The leftmost large branch}
Let \(M\) satisfy the induction scheme and the least-number principle, let
\(\phi\), \(\theta\), \(E\), \(L\), \(S\) and \(\sigma\) be as in
\cref{def:p03-w03-p03-030-tree-search}, and let \(\theta(M)\) be unbounded in
\(M\).  Then:
\begin{enumerate}
\item for every rank \(\bar u\), \(L(\bar u,\cdot,\cdot)\) induces a strict
  linear order on the \(E(\bar u,\cdot,\cdot)\)-classes;
\item for every rank \(\bar u\), \(S(\bar u,\cdot)\) defines exactly one
  \(E(\bar u,\cdot,\cdot)\)-class, and that class is large;
\item if \(\bar w<_{n+1}\bar u\) and \(S(\bar u,y)\), then \(S(\bar w,y)\);
\item if \(\bar x<_{n+1}\bar u\) and \(S(\bar u,y)\), then
  \(\phi(\bar x,y)\leftrightarrow\sigma(\bar x)\).
\end{enumerate}
\end{hblemma}

\begin{hbtheorem}{thm:p03-w03-p03-030-mss}{proved=annex:p03-w03-p03-030-mss}{Self-extending theories from the induction scheme}
Let \(T\) be a consistent theory in a countable first-order vocabulary \(L\)
containing \(<\).  Suppose that \(T\) contains the induction scheme and the
least-number principle and is a Skolem theory for first-order formulas.
Then to each first-order formula \(\phi(x_0,\ldots,x_{n-1},y)\) there is
associated a first-order formula \(\sigma_\phi(x_0,\ldots,x_{n-1})\) such that
\(T\) can be extended to a self-extending theory \(T'\) containing the
universal closures of all formulas
\begin{gather*}
 Qy\,\phi(\bar x,y)\leftrightarrow\sigma_\phi(\bar x)
 \qquad\text{and}\\
 \forall z\,\bigl(Qx\,\exists y\,(\phi(x,y)\wedge y<z)\rightarrow
 \exists y\,Qx\,\phi(x,y)\bigr),
\end{gather*}
and also the sentence \(\forall y\,Qx\,(y<x)\).  Moreover, every model \(M\)
of \(T\) expands to a reduced model \((M,q_M)\) of \(T'\).
\end{hbtheorem}

This is Theorem~6.2.2 of the source
\cite[p.~206]{Schmerl1985TransferTheorems}, whose proof is on
\cite[pp.~207--208]{Schmerl1985TransferTheorems}, with two differences of
hypothesis.  The source assumes only the induction scheme.  It remarks that
adjoining all definable terms then yields a Skolem theory, and it treats
such theories as Skolem theories from the outset
\cite[p.~206]{Schmerl1985TransferTheorems}.  The proof written here also
uses the least-number principle, to find the first parameter tuple at which
two classes differ and to find least witnesses.  Over the induction scheme
the least-number principle, a scheme, is equivalent to one sentence: every
element has an immediate successor (\cref{thm:p03-w03-p03-030-is-successor}).
Both principles hold in every structure whose order has type \(\omega\),
since there every nonempty set has a least element and every bounded set is
finite.  The theorem for Skolem theories that satisfy the induction scheme
but not the least-number principle is not proved here.  Such theories exist,
and for some of them adjoining the functions definable without parameters
does not produce a Skolem theory
(\cref{def:p03-w03-p03-030-example-one-plus-z}).

\begin{hblemma}{thm:p03-w03-p03-030-is-successor}{proved=annex:p03-w03-p03-030-is-successor}{The least-number principle over the induction scheme}
Let \(M\) satisfy the induction scheme of
\cref{def:p03-w03-p03-030-induction-scheme} for all formulas with parameters
in its vocabulary, and let \(0\) be its first element.
\begin{enumerate}
\item Every nonempty definable subset of \(M\) with an upper bound has a
  greatest element, and every element other than \(0\) has an immediate
  predecessor.
\item \(M\) satisfies the least-number principle if and only if
  \[
   M\entails\forall x\,\exists y\,\bigl(x<y\wedge\forall z\,(x<z\rightarrow
   y\le z)\bigr),
  \]
  that is, if and only if every element has an immediate successor.
\end{enumerate}
\end{hblemma}

\begin{hbexample}{def:p03-w03-p03-030-example-one-plus-z}{The induction scheme without the least-number principle}
Let \(\mathfrak Z\) be the linear order whose domain is \(\{e\}\cup\mathbb
Z\), where \(e\) is a new point placed below every integer and the integers
keep their usual order.  Then \(e\) is the first element and there is no last
element.  Every nonempty subset \(X\) with an upper bound has a greatest
element.  If \(X=\{e\}\), that element is \(e\).  Otherwise \(X\cap\mathbb Z\)
is a nonempty set of integers, bounded above by the upper bound of \(X\),
which is then an integer; its greatest integer is the greatest element of
\(X\).  So every nonempty subset without a greatest element is unbounded, that
is, cofinal.  Hence the induction scheme holds in \(\mathfrak Z\) and in every
expansion of \(\mathfrak Z\) by further relations, functions and constants,
for all formulas of the expanded vocabulary.  The least-number principle
fails: the formula \(\exists u\,(u<x)\) defines \(\mathbb Z\), which has no
least element.  As \cref{thm:p03-w03-p03-030-is-successor} predicts, \(e\)
has no immediate successor.

\emph{Skolem expansions exist.}  List the domain as
\(e,0,1,-1,2,-2,\ldots\).  Let \(L_0=\{<\}\), and let \(L_{k+1}\) add to
\(L_k\) a function symbol \(g_\chi\) for each \(L_k\)-formula
\(\chi(\bar z,y)\), interpreted as the first element of the list satisfying
\(\chi(\bar z,y)\) when there is one and as \(e\) otherwise.  The expansion to
\(\bigcup_kL_k\) is a Skolem structure: every formula lies in some \(L_k\),
and \(g_\chi\) picks a witness whenever one exists.  By the previous
paragraph its theory contains the induction scheme, and it does not contain
the least-number principle.  So it is a theory to which the source's
Theorem~6.2.2 applies and \cref{thm:p03-w03-p03-030-mss} does not.

\emph{Definable functions do not suffice.}  The map \(\pi\) that fixes \(e\)
and sends each integer \(k\) to \(k+1\) is an automorphism of
\(\mathfrak Z\).  Every relation or function definable in \(\mathfrak Z\)
without parameters is preserved by \(\pi\), because an automorphism preserves
the satisfaction of every formula without parameters.  So \(\pi\) is an
automorphism of every expansion of \(\mathfrak Z\) by such relations and
functions.  Repeating the expansion does not help: what is definable without
parameters in an expansion that \(\pi\) preserves is again preserved by
\(\pi\).  By induction on terms, every closed term of such an expansion
denotes a fixed point of \(\pi\), and the only fixed point is \(e\).  The
formula \(\exists u\,(u<y)\) is satisfied by every integer and not by
\(e\).  So no closed term is a witness for it, and the expansion is not
a Skolem structure.  The passage from the induction scheme to a Skolem
theory by adjoining definable terms therefore needs more than the induction
scheme when the terms are taken without parameters.  The source's remark is
not checked here under any other reading.
\end{hbexample}

\begin{hbcorollary}{thm:p03-w03-p03-030-end-extension-pa}{proved=annex:p03-w03-p03-030-end-extension-pa}{Proper elementary end-extensions}
Let \(M\) be a structure for a countable vocabulary containing \(<\) that
satisfies the induction scheme and the least-number principle.  Then \(M\)
has a proper elementary end-extension \(N\) with \(\abs N=\abs M\).
\end{hbcorollary}

The source states this consequence for every model of the induction scheme
in a countable vocabulary \cite[p.~207]{Schmerl1985TransferTheorems}; the
version proved here adds the least-number principle, as in
\cref{thm:p03-w03-p03-030-mss}, or equivalently the sentence that every
element has an immediate successor
(\cref{thm:p03-w03-p03-030-is-successor}).  For Peano arithmetic both schemes
follow from the induction axioms, so the corollary yields MacDowell and Specker's
theorem.

\begin{hbcorollary}{thm:p03-w03-p03-030-pa}{proved=annex:p03-w03-p03-030-pa}{The MacDowell--Specker theorem for Peano arithmetic}
Let \(\mathcal M\) be a model of Peano arithmetic \(\PA\), that is, of the
seven axioms of Robinson arithmetic \(\mathsf Q\) together with the
induction scheme for all formulas with parameters, in the language
\(\Lang_A=\{0,S,+,\times\}\) of \cref{def:language-arithmetic}, and order
it by \(x<y\) iff \(\exists z\,(x+Sz=y)\).  Then \(\mathcal M\) has a proper
elementary extension \(\mathcal N\) of the same cardinality such that every
element of \(\mathcal N\) that is \(<\)-below an element of \(\mathcal M\)
belongs to \(\mathcal M\).
\end{hbcorollary}

The annex proof derives, from the axioms of \(\mathsf Q\) and the induction
scheme, that \(<\) is a linear order with first element \(0\) and no last
element, that every element \(x\) has the immediate successor \(Sx\), and
that the induction scheme of \cref{def:p03-w03-p03-030-induction-scheme} and
the least-number principle hold for every formula of \(\Lang_A\cup\{<\}\).
It then applies \cref{thm:p03-w03-p03-030-end-extension-pa} to the expansion
of \(\mathcal M\) by the defined order.  The source states only that the
MacDowell--Specker theorem concerns models of Peano arithmetic, without this
verification \cite[p.~206]{Schmerl1985TransferTheorems}.

\begin{hbcorollary}{thm:p03-w03-p03-030-like-transfer}{proved=annex:p03-w03-p03-030-like-transfer}{From \(\omega\)-like models to \(\lambda\)-like models}
Let \(\Sigma\) be a countable set of first-order sentences in a countable
vocabulary containing \(<\).  Suppose every finite subset of \(\Sigma\) has
a model whose order \(<\) has order type \(\omega\).  Then for every
uncountable cardinal \(\lambda\), \(\Sigma\) has a \(\lambda\)-like model.
\end{hbcorollary}

The corollary is the case \(\mu=\aleph_0\) of Schmerl's Theorem~6.2.1.  To
state that theorem, recall the source's transfer notation for cardinal-like
models: for infinite cardinals \(\kappa\), \(\lambda\) and \(\mu\),
\(\kappa\to\lambda\) \emph{\(\mu\)-compactly} means that every set of at most
\(\mu\) first-order sentences, each finite subset of which has a
\(\kappa\)-like model, has a \(\lambda\)-like model
\cite[Definition~1.3.1, p.~180]{Schmerl1985TransferTheorems}.  Up to
isomorphism, \((\omega,<)\) is the only \(\aleph_0\)-like linear order
\cite[p.~179]{Schmerl1985TransferTheorems}, so an \(\aleph_0\)-like model is
a model whose order has type \(\omega\).

\begin{hbtheorem}{thm:p03-w03-p03-030-transfer-mu}{proved=annex:p03-w03-p03-030-transfer-mu}{Transfer from \(\aleph_0\)-like to \(\lambda\)-like models}
If \(\aleph_0\le\mu<\lambda\), then \(\aleph_0\to\lambda\)
\(\mu\)-compactly.  That is, let \(\Sigma\) be a set of at most \(\mu\)
first-order sentences in a vocabulary containing \(<\), each finite subset of
which has a model whose order has type \(\omega\).  Then \(\Sigma\) has a
\(\lambda\)-like model.
\end{hbtheorem}

This is Theorem~6.2.1 of the source
\cite[p.~206]{Schmerl1985TransferTheorems}, stated there without proof.  The
source says that it follows from Theorem~6.2.2 and attributes to Fuhrken the
observation that it follows from MacDowell and Specker's theorem
\cite[pp.~206--207]{Schmerl1985TransferTheorems}.  The proof in the annex is
written for this handbook and takes a different route, which does not need
\cref{thm:p03-w03-p03-030-mss} and works for vocabularies of any size
\(\mu\).  It fixes one nonprincipal ultrafilter on \(\mathbb N\), expands
the finite approximations, all of which live on \(\mathbb N\), by
least-witness functions, and obtains by
\cref{thm:p03-w03-p03-030-compactness} a self-extending theory of size \(\mu\)
that contains \(\Sigma\) and the bounded scheme (B).  The iteration of
\cref{sec:p03-w03-p03-030-iteration} then runs for \(\lambda\) steps.  The
hypothesis \(\mu<\lambda\) is what keeps every stage of the chain below
\(\lambda\).

\begin{hbexample}{def:p03-w03-p03-030-example-mu-lambda}{The condition \(\mu<\lambda\) cannot be weakened to \(\mu\le\lambda\)}
Let \(\lambda\) be an infinite cardinal.  Take constants \(c_\alpha\)
\((\alpha<\lambda)\) and \(d\), and let \(\Sigma\) say that \(<\) is a linear
order, that \(c_\alpha<d\) for every \(\alpha\), and that
\(c_\alpha\ne c_\beta\) for \(\alpha\ne\beta\).  Then \(\abs\Sigma=\lambda\).
A finite subset mentions finitely many constants, and it holds in
\((\omega,<)\) when those \(c_\alpha\) are sent to distinct numbers and \(d\)
to a larger one.  But in any model of \(\Sigma\) the element \(d\) has at
least \(\lambda\) elements below it, so no model of \(\Sigma\) is
\(\lambda\)-like.  So \(\aleph_0\to\lambda\) fails \(\lambda\)-compactly.  With
\(\lambda=\aleph_0\) this also shows that the theorem has no analogue with
\(\lambda=\mu=\aleph_0\).
\end{hbexample}

\section{Limits of the method and of this presentation}
\label{sec:p03-w03-p03-030-limits}

\begin{hbwarning}{def:p03-w03-p03-030-limits}{What the hypotheses do and what is not established here}
\begin{enumerate}
\item The Skolem property and the ultrafilter schemes are both needed for the
  canonical extension to be elementary.  Even with both, a self-extending
  theory need not give a least new element
  (\cref{def:p03-w03-p03-030-example-n}); the strongly \(\lambda\)-like
  conclusion needs the regressive scheme together with
  \(\forall y\,Qx\,(y<x)\).
\item The chain of \cref{thm:p03-w03-p03-030-strongly-like} has length a
  regular \(\lambda\) larger than the language.  Regularity is used to keep
  every stage below \(\lambda\) at limit stages, and the notion of a
  strongly \(\lambda\)-like order is defined in the source only for regular
  uncountable \(\lambda\).
\item The proof of \cref{thm:p03-w03-p03-030-mss} lists all formulas in a
  sequence of length \(\omega\) and builds the decreasing sequence
  \(\theta_0,\theta_1,\ldots\) along it.  It therefore uses the countability
  of the vocabulary.  The proof of \cref{thm:p03-w03-p03-030-transfer-mu}
  avoids this, because all the finite approximations there have order type
  \(\omega\); it does not give end-extensions of arbitrary models in an
  uncountable vocabulary.
\item In the formula the source gives for \(\theta_n\), the last disjunct
  reads \(y\le z\) \cite[p.~207]{Schmerl1985TransferTheorems}.  With the
  outer condition \(z<y\), substituting \(w\defeq y\) makes the formula
  unsatisfiable, so the set it defines would be empty.  The annex uses
  \(y\le w\), which makes \(y\) the least suitable element above \(z\).
\item The source also states, without proof, that a theory with suitable
  two-cardinal models, with \(\kappa\)-like models for \(n\)-Mahlo
  \(\kappa\), or with \((\kappa,X)\)-like models for \(n\)-subtle \(X\) for
  every \(n<\omega\), extends to a self-extending theory containing the
  bounded or regressive scheme.  For the \(n\)-subtle case it sketches the
  proof and derives from it a transfer to strongly \(\kappa\)-like models
  \cite[Theorems~6.1.2, 6.1.3, 6.1.5 and 6.1.6, pp.~204--206]{Schmerl1985TransferTheorems}.
  In its final remarks it states, again without proof, infinitary versions:
  a transfer theorem for \(\mathscr L_{\omega_1\omega}\) with a cardinality
  quantifier, the transfer to \(\lambda\)-like and strongly
  \(\lambda\)-like models for a sentence of an admissible fragment under
  Mahlo and subtle hypotheses indexed below the well-ordering number of the
  fragment, and the Hanf numbers of admissible fragments
  \cite[Theorems~7.2.1--7.2.3, p.~209]{Schmerl1985TransferTheorems}.  None of
  these is proved in this chapter, and none of them is used in a proof here.
\end{enumerate}
\end{hbwarning}

\paragraph{Sources.}  The definitions, the canonical extension, the
iteration and the tree argument follow
\cite[\S6, pp.~202--208]{Schmerl1985TransferTheorems}, from the 1985 edition
of \emph{Model-Theoretic Logics} \cite{BarwiseFeferman1985}.  The
indiscernibility theorem, the hull, union and compactness lemmas, the
examples, the verification for Peano arithmetic, the proof of the transfer
theorem from \(\aleph_0\)-like to \(\lambda\)-like models and the full
verifications in the annex are written for this handbook.  The statement of
that transfer theorem and the transfer notation are the source's
\cite[Definition~1.3.1, p.~180; Theorem~6.2.1, p.~206]{Schmerl1985TransferTheorems}.
